\section{Markov chains and boundary laws}

Throughout, $S$ is the vertex set of a locally finite graph $G$ (a tree where stated), $E$ a
countable set with the discrete $\sigma$-algebra and counting measure $\lambda$, and for
$\Lambda \in \mathcal S$, $\partial\Lambda$ its outer boundary and $\operatorname{bonds}(\Lambda)$
the edges meeting $\Lambda$.

\begin{definition}[Markov specification, Markov chain on a graph, Georgii (12.1)--(12.2)]
    \label{def:tr-markov}
    \uses{def:spec, def:cylinder-event}
    \lean{MeasureTheory.GibbsMeasure.Tree.IsMarkovSpecification, MeasureTheory.GibbsMeasure.Tree.IsMarkovChain, SimpleGraph.past}
    \leanok

    A specification $\gamma$ is \textbf{Markov} if $\gamma_\Lambda(\sigma_\Lambda = \zeta_\Lambda \mid \cdot)$
    is $\mathcal F_{\partial\Lambda}$-measurable for every $\Lambda$. A probability measure $\mu$ is
    a \textbf{Markov chain} on $G$ if for every edge $ij$ and $y \in E$,
    $\mu(\sigma_j = y \mid \mathcal F_{\operatorname{past}(i,j)}) = \mu(\sigma_j = y \mid \sigma_i)$,
    where $\operatorname{past}(i,j)$ is the set of sites closer to $i$ than to $j$.
\end{definition}

\begin{definition}[Transfer families and $\gamma^Q$, Georgii (12.7)--(12.9)]
    \label{def:tr-transfer}
    \uses{def:tr-markov, def:premodif}
    \lean{MeasureTheory.GibbsMeasure.Tree.IsTransferFamily, MeasureTheory.GibbsMeasure.Tree.transferWeight, MeasureTheory.GibbsMeasure.Tree.transferSpecification, MeasureTheory.GibbsMeasure.Tree.isMarkovSpecification_transferSpecification}
    \leanok

    A \textbf{transfer family} is a symmetric family $Q = (Q_{ij})$ of positive matrices with
    finite entries on the edges of $G$, whose partition functions $Z_\Lambda$ are finite. The
    specification $\gamma^Q$ is the $\lambda$-specification of the weights
    $\prod_{b \in \operatorname{bonds}(\Lambda)} Q_b(\sigma)$; it is Markov.
\end{definition}

\begin{definition}[Boundary law on a graph, Georgii (12.10)]
    \label{def:tr-boundary-law}
    \uses{def:tr-transfer}
    \lean{MeasureTheory.GibbsMeasure.Tree.IsBoundaryLaw, MeasureTheory.GibbsMeasure.Tree.Markov.IsBoundaryLaw.isBoundaryLaw_hasse_int}
    \leanok

    A \textbf{boundary law} for $Q$ is a family $(\ell_{ij})$ of positive finite functions on $E$,
    indexed by the oriented edges, with
    $\ell_{ij} \propto \prod_{k \in \partial i \setminus \set j} \ell_{ki} Q_{ki}$ for every
    oriented edge $ij$, and finite mass $\sum_x \prod_{k \in \partial i}(\ell_{ki}Q_{ki})(x)$.
    On $\mathbb Z$ a boundary law in the sense of Chapter 11 is a boundary law on the Hasse graph
    (Georgii (12.11)).
\end{definition}

\begin{theorem}[Georgii (12.12)]
    \label{thm:tr-12-12}
    \uses{def:tr-boundary-law, def:gibbs-meas, thm:bl-11-9a}
    \lean{MeasureTheory.GibbsMeasure.Tree.boundaryLawMeasure, MeasureTheory.GibbsMeasure.Tree.IsBoundaryLaw.boundaryLawMeasure_cyl, MeasureTheory.GibbsMeasure.Tree.IsBoundaryLaw.isGibbsMeasure_transferSpecification_boundaryLawMeasure, MeasureTheory.GibbsMeasure.Tree.IsBoundaryLaw.isMarkovChain_boundaryLawMeasure, MeasureTheory.GibbsMeasure.Tree.IsMarkovChain.exists_isBoundaryLaw_eq_boundaryLawMeasure}
    \leanok

    Let $G$ be a tree and $Q$ a transfer family. (a) Every boundary law $\ell$ defines a
    probability measure $\mu_\ell$ on $\Cfg$ with
    \[
        \mu_\ell(\sigma_{\Lambda \cup \partial\Lambda} = \zeta)
        = z_\Lambda \prod_{k \in \partial\Lambda} \ell_{k k_\Lambda}(\zeta_k)
          \prod_{b \in \operatorname{bonds}(\Lambda)} Q_b(\zeta) \tag{12.13}
    \]
    for connected $\Lambda$, and $\mu_\ell \in \mathcal G(\gamma^Q)$ is a Markov chain. (b)
    Conversely every Markov chain in $\mathcal G(\gamma^Q)$ is $\mu_\ell$ for some boundary law
    $\ell$.
\end{theorem}
\begin{proof}
    \leanok
    (a) The right-hand side of (12.13) is consistent along connected volumes, by induction over
    the leaves attached at the boundary (12.14); Kolmogorov extension; the DLR equation for
    singletons follows from (1.33). (b) $\ell_{ij}(x) = P_{ji}(a,x)/Q_{ji}(a,x)$ for the
    transition probabilities $P$ of the chain and a reference $a$; the Markov property and the
    Gibbs property identify the cylinder probabilities with (12.13).
\end{proof}

\begin{proposition}[Georgii (12.15)--(12.17)]
    \label{prop:tr-12-16}
    \uses{def:tr-boundary-law}
    \lean{MeasureTheory.GibbsMeasure.Tree.IsBoundaryLaw.eq_prod_div_of_normalized, MeasureTheory.GibbsMeasure.Tree.isBoundaryLaw_const_iff}
    \leanok

    Normalised boundary laws ($\ell_{ij}(a) = 1$) satisfy $\ell_{ij} = \prod_{k \in \partial i \setminus\set j}(\ell_{ki}Q_{ki})/(\ell_{ki}Q_{ki})(a)$.
    On a regular tree of degree $d + 1$ with a constant transfer matrix $Q_0$, a constant
    normalised $\ell_0$ is a boundary law iff $\ell_0(x) = \bigl(\ell_0 Q_0(x)/\ell_0 Q_0(a)\bigr)^d$
    and $\sum_x (\ell_0 Q_0(x))^{d+1} < \infty$.
\end{proposition}

\begin{proposition}[Georgii (12.3)(2),(4), (12.5); uniqueness in (12.12)(b)]
    \label{prop:tr-12-3}
    \uses{def:tr-markov, thm:tr-12-12}
    \lean{MeasureTheory.GibbsMeasure.Tree.IsMarkovChain.exists_parent_measure_cyl_eq, MeasureTheory.GibbsMeasure.Tree.transitionProb_mul_transitionProb_swap_eq, MeasureTheory.GibbsMeasure.Tree.tsum_transitionProb_eq_one, MeasureTheory.GibbsMeasure.Tree.IsBoundaryLaw.exists_const_mul_eq_of_boundaryLawMeasure_eq}
    \leanok

    For a Markov chain $\mu$ on a tree, the cylinder probabilities factorise along the edges
    from a root (12.4), the transition probabilities $P_{ij}$ are stochastic and satisfy the
    reversibility relation (12.5). Two boundary laws whose measures (12.13) coincide differ by a
    positive constant on each oriented edge.
\end{proposition}

\begin{theorem}[Georgii (12.17), the correspondence]
    \label{thm:tr-12-17}
    \uses{prop:tr-12-16, prop:tr-12-3}
    \lean{MeasureTheory.GibbsMeasure.Tree.IsCompletelyHomogeneousMarkovChain, MeasureTheory.GibbsMeasure.Tree.isCompletelyHomogeneousMarkovChain_iff_exists_isBoundaryLaw_const_solves, MeasureTheory.GibbsMeasure.Tree.eq_of_isBoundaryLaw_const_boundaryLawMeasure_eq}
    \leanok

    On a regular tree with constant transfer matrix $Q_0$, $\mu \in \mathcal G(\gamma^{Q_0})$ is
    a completely homogeneous Markov chain iff $\mu$ is the measure of a constant normalised
    boundary law solving (12.16); that solution is unique.
\end{theorem}

\begin{theorem}[Georgii (12.6)]
    \label{thm:tr-12-6}
    \uses{def:tr-markov, thm:ext-77-concrete}
    \lean{MeasureTheory.GibbsMeasure.Tree.IsMarkovSpecification.dependsOn_apply, MeasureTheory.GibbsMeasure.Tree.exists_isMarkovChain_of_mem_extremePoints}
    \leanok

    Let $G$ be a locally finite tree with countable vertex set and $\gamma$ a Markov
    specification on $E^S$, $E$ countable. Then every extreme $\mu \in \mathcal G(\gamma)$ is a
    Markov chain on $G$.
\end{theorem}
\begin{proof}
    \leanok
    Fix an oriented edge $ij$ and exhaust the past of $ij$ by connected volumes $\Lambda(n)$.
    By the Markov property, $\mu(\sigma_j = y \mid \mathcal F_{\Lambda(n)^c})$ depends only on
    $\sigma_i$ and the configuration outside $\Lambda(n)$; the reversed martingale converges
    to $\mu(\sigma_j = y \mid \mathcal T_{\mathrm{past}})$, which by tail triviality (7.7)(a) is
    $\mu(\sigma_j = y \mid \sigma_i)$ after freezing the single coordinate $\sigma_i$ ($E$ countable).
\end{proof}

\begin{corollary}[Georgii (12.18)]
    \label{cor:tr-12-18}
    \uses{thm:tr-12-12, prop:tr-12-3}
    \lean{MeasureTheory.GibbsMeasure.Tree.not_isMarkovChain_sum_smul_of_forall_adj_exists_boundaryLaw_eq}
    \leanok

    Let $\mu_1, \dots, \mu_N$ be distinct measures of normalised boundary laws $\ell^{(n)}$ for
    a transfer family on a locally finite tree such that every oriented edge $ji$ has a neighbour
    $k \neq j$ of $i$ with $\ell^{(n)}_{ki} = \ell^{(n)}_{ji}$ for all $n$. Then no nontrivial
    convex combination of the $\mu_n$ is a Markov chain.
\end{corollary}
\begin{proof}
    \leanok
    A convex combination is Gibbs; if it were a Markov chain it would be the measure of a
    boundary law $\ell$ by Theorem \ref{thm:tr-12-12}. The cylinder formula (12.13) at
    singleton volumes gives $\ell_{ji}(x)\ell_{ji}(y) = \sum_n w_n \ell^{(n)}_{ji}(x)\ell^{(n)}_{ji}(y)$
    and $\ell_{ji}(z) = \sum_n w_n \ell^{(n)}_{ji}(z)$; the equality case of Jensen's inequality
    for $x \mapsto x^2$ forces all $\ell^{(n)}_{ji}$ with $w_n > 0$ to agree, hence the
    corresponding $\mu_n$, contradicting distinctness.
\end{proof}
